%%%% RESUMO

\begin{resumoutfpr}%% Ambiente resumoutfpr
O transporte público é um componente vital para a mobilidade em cidades inteligentes, gerando grandes volumes de dados de geolocalização que podem ser explorados para aprimorar a gestão operacional. No entanto, a utilização direta de dados brutos de GPS enfrenta desafios significativos devido a ruídos, imprecisões e falhas de comunicação, dificultando a reconstrução fiel dos itinerários. Este trabalho propõe uma metodologia para aplicar técnicas de Mineração de Processos aos dados da Rede Integrada de Transporte (RIT) de Curitiba (URBS). O objetivo principal é transformar registros desestruturados de GPS em um log de eventos confiável, permitindo a descoberta e análise dos processos reais de operação. A metodologia adota a biblioteca PM4Py para a mineração e integra algoritmos de Map Matching para a identificação precisa de paradas em pontos de ônibus, além de técnicas de interpolação temporal para correção de lacunas nos dados. Como resultado, foi gerado um arquivo XES estruturado que possibilitou a aplicação de algoritmos de descoberta (Inductive Miner) e verificação de conformidade. A análise permitiu identificar gargalos de desempenho e desvios entre as rotas executadas e as planejadas, demonstrando a viabilidade da engenharia de dados espaciais como pré-requisito para a inteligência em transporte.
\end{resumoutfpr}
